\chapter*{ВВЕДЕНИЕ}
\addcontentsline{toc}{chapter}{ВВЕДЕНИЕ}

В условиях активного развития цифровых технологий возрастает роль информационных систем, обеспечивающих наглядное представление сложных данных. Особенно это актуально для информации, обладающей одновременно пространственной и временной привязкой. К таким данным относятся сведения об исторических военных конфликтах, включающие информацию о сторонах противостояния, операциях, перемещении армейских соединений, изменении линии фронта и связанных событиях. Традиционные способы представления подобных данных в виде текстовых описаний, таблиц и статических карт не позволяют в полной мере отразить динамику произошедших событий и их масштаб.

Актуальность темы обусловлена необходимостью создания программного средства, которое обеспечивало бы хранение и визуализацию данных об исторических военных действиях в единой системе. Существующие решения, как правило, либо ориентированы только на отображение отдельных картографических данных, либо не обеспечивают временной навигации, либо не предоставляют единой среды, объединяющей базу данных, серверную логику, пользовательский интерфейс и средства визуального анализа. В связи с этим разработка веб-приложения, позволяющего обрабатывать пространственно-временные данные об исторических военных событиях, является практически значимой и востребованной задачей.

Цель работы заключается в разработке веб-приложения «Палантир», предназначенного для хранения, обработки и визуализации данных об исторических военных событиях.

Для достижения поставленной цели необходимо решить следующие задачи:
\begin{enumerate}
\item Провести анализ предметной области исторических военных конфликтов, выполнить обзор аналогичных программных решений и определить требования к веб-приложению визуализации данных об исторических военных событиях.
\item Определить основные технологические и архитектурные решения, необходимые для реализации веб-приложения.
\item Разработать базу данных и программную реализацию веб-приложения «Палантир», включающие основные функции, связанные с хранением, обработкой и визуализацией данных.
\item Провести тестирование разработанной системы и рассмотреть подготовку приложения к развёртыванию.
\end{enumerate}

Практическая значимость работы заключается в создании веб-информационной системы, которая может служить основой для визуального представления данных об исторических военных событиях, а также для дальнейшего расширения функциональности и развития проекта.
